\documentclass[11pt]{article}
% Shared preamble for all solution files.
\usepackage[margin=1.1in]{geometry}
\usepackage{enumitem}
\usepackage{tikz}
\usepackage[hidelinks]{hyperref}
\usepackage{amsmath,amssymb,amsthm}
\newcommand{\fa}{\mathfrak a}
\newcommand{\fb}{\mathfrak b}
\newcommand{\fc}{\mathfrak c}
\newcommand{\fp}{\mathfrak p}
\newcommand{\fq}{\mathfrak q}
\newcommand{\fm}{\mathfrak m}
\newcommand{\fn}{\mathfrak n}
\newcommand{\fr}{\mathfrak r}
\newcommand{\Nil}{\mathfrak N}
\newcommand{\Jac}{\mathfrak R}
\newcommand{\rad}{r}
\newcommand{\Spec}{\operatorname{Spec}}
\newcommand{\Max}{\operatorname{Max}}
\newcommand{\Ann}{\operatorname{Ann}}
\newcommand{\Ker}{\operatorname{Ker}}
\newcommand{\Z}{\mathbf Z}
\newcommand{\Q}{\mathbf Q}
\newcommand{\R}{\mathbf R}
\newcommand{\C}{\mathbf C}
\newcommand{\ci}{\mathrm{c}}
\newcommand{\ex}{\mathrm{e}}

\setlength{\parindent}{0pt}
\setlength{\parskip}{0.55em}
% \exercise{label}{statement}  -- followed by the solution text
\newcommand{\exercise}[2]{\par\bigskip\noindent{\large\textbf{Exercise #1}}\par\smallskip
  \begin{quote}\small\textbf{Problem.}\ #2\end{quote}\par}
\newcommand{\pf}{\noindent\textbf{Solution.}\ }

\usepackage{needspace}
\AddToHook{cmd/section/before}{\Needspace{14\baselineskip}}
\let\chapterexercise\exercise
\renewcommand{\exercise}[2]{\Needspace{8\baselineskip}\chapterexercise{#1}{#2}}
\title{Liu, Chapter 9: Regular surfaces\\[2pt]\large Solutions and corrections to the exercises}
\author{}\date{}
\begin{document}\maketitle
\textbf{Conventions.} These 54 exercises follow the four lists on printed pp.~391--394, 408--410, 426--429 and 449--453 of the supplied edition. An arithmetic surface is regular, integral, projective and flat over a one-dimensional Dedekind scheme. Intersections include residue-field degrees; $K$ denotes a relative canonical divisor. We use the factorization, adjunction, contraction and minimal-model theorems of this chapter. An exceptional divisor means a contractible curve with regular image, not an arbitrary component of an exceptional locus. Unless specified, statements about a single fibre are made after localization at its base point.

\textbf{Corrections.} Some printed assertions are false with their stated hypotheses. We impose a common regular point to obtain the dimension equality in 1.2, and distinguish tangent spaces from conormal lines in 1.3, total degree from multidegree in 1.12, and a partial fibre from the full primitive fibre in 1.14. Exercise 2.6 requires control of the remaining fibre singularities; the numerical bound in 2.11(d) needs correction. In 3.5 a global finite etale refinement need not exist; 3.8 must keep track of the generic constant field, and 3.9(b) needs minimality. The direct-sum assertions in 4.3(b) omit gluing terms. The algorithmic assertion in 4.7(c) need not recover a prescribed resolution. Exercise 4.8 has a local answer, but no uniform global answer without further hypotheses. The effectivity and triviality assertions in 4.16 are local on the Dedekind base. Counterexamples and corrected statements are supplied at the corresponding labels; a counterexample is a resolution of a false question, rather than a proof of its printed assertion.

\section*{9.1 Intersection theory on a regular surface}

\exercise{1.1}{Move a Cartier divisor on a quasi-projective variety over an infinite field away from finitely many prescribed nongeneric points, without using Proposition 1.11.}
\pf Choose an ample invertible sheaf $A$. For $n$ large, both $A^n$ and $\mathcal O(D)\otimes A^n$ are generated by finitely many global sections. For each prescribed point, the sections vanishing in its fibre form a proper linear subspace of the finite-dimensional space of generators. A finite union of proper subspaces cannot exhaust a vector space over an infinite field. Choose sections $t$ and $s$ of the two sheaves avoiding every point, and also every associated point of $X$, so that their zero schemes are effective Cartier divisors. Then $D'=\operatorname{div}(s)-\operatorname{div}(t)$ represents $D$ and avoids all the prescribed points. On an integral variety the associated-point condition merely requires the sections to be nonzero. \hfill$\square$

\exercise{1.2}{Move $n=\dim X$ Cartier divisors so that the successive intersections of their supports have dimension $n-r$.}
\pf Fix a closed point $x$ in the nonempty regular locus of $X$, and a basis $t_1,\ldots,t_n$ of $\fm_x/\fm_x^2$. We construct $D_i'$ through $x$, with local equation congruent to $t_i$ modulo $\fm_x^2$, while avoiding the generic points of the preceding support intersection.

Here is why these requirements can be imposed simultaneously. For an ample $A$ and sufficiently large $N$, write $D_i'=\operatorname{div}(s)-\operatorname{div}(u)$ with $s$ a section of $\mathcal O(D_i)\otimes A^N$ and $u$ a section of $A^N$. Serre's theorem supplies arbitrary first jets at $x$ and generation away from $x$ even after those jets are fixed. Choose $u$ nonvanishing at $x$ and at the finitely many prescribed generic points. Choose the jet of $s/u$ to be $t_i$, and choose $s$ nonvanishing at those generic points. Each forbidden vanishing condition is a proper affine linear subspace of the space with prescribed jet; an infinite field permits avoiding all of them. Both sections can also be chosen nonzero at the generic point of $X$.

Inductively no component of the preceding intersection is contained in the next support, so its dimension is at most $n-r$. At $x$ the first $r$ local equations are part of a regular system of parameters, so their common zero set has local dimension $n-r$. Since all denominators are units at $x$, it lies in the intersection of the supports. This gives the reverse inequality and proves the requested equality. For $r<n$, the preceding intersection has positive-dimensional local components at $x$, so the generic points being avoided are different from $x$; there is no incompatibility in the choices. \hfill$\square$

\exercise{1.3}{Characterize normal crossings by regularity of the prime components and linear conditions on their tangent spaces.}
\pf \textbf{Replace the printed direct sum of tangent spaces by independence of conormal lines.} Only components through $x$ are relevant. Put $A=\mathcal O_{X,x}$, with maximal ideal $\fm$. If $f_i$ generates the ideal of such a component, the correct criterion is that the classes $\overline f_i\in\fm/\fm^2$ be linearly independent. These classes then extend to a basis, hence the $f_i$ extend to a regular system of parameters. The local equation of $D$ is a unit times $\prod f_i^{n_i}$, exactly the normal-crossings condition. Conversely such parameters give independent classes and regular components. Equivalently, the lines $k(x)\,df_i\subseteq T_{X,x}^{\vee}$ are independent. In $\mathbf A^3$, the divisor $uv=0$ has normal crossings, while its two tangent planes have a common line and cannot form a direct sum. If normal crossings is defined after an etale extension, apply this criterion to the resulting local branches; the criterion with the original prime components describes strict normal crossings. \hfill$\square$

\exercise{1.4}{Show that properties (a) and (c) alone uniquely determine the bilinear intersection pairing in Theorem 1.12.}
\pf Fix a vertical prime $\Gamma$. By the valuation moving lemma choose $D'\sim D$ with no component $\Gamma$. Property (c) identifies $i_s(D,\Gamma)$ with $i_s(D',\Gamma)$, and (a) specifies the latter as the sum of the local intersections with residue degrees. Bilinearity then determines the value for every vertical divisor $\sum n_i\Gamma_i$. Symmetry need not be assumed for this uniqueness argument. \hfill$\square$

\exercise{1.5}{For two distinct smooth plane curves through a rational point $x$, show that $i_x(D,E)\ge2$ exactly when their tangent lines coincide.}
\pf In the regular local ring $A=\mathcal O_{\mathbf P^2,x}$ let $f,g$ be their equations. Both have nonzero linear terms. The quotient $A/(f,g)$ has length one exactly when $(f,g)=\fm$, which by Nakayama holds exactly when these linear terms are independent. In the two-dimensional cotangent space, dependence is equivalent to equality of the tangent lines. Otherwise the finite positive length is at least two. \hfill$\square$

\exercise{1.6}{For a vertical rational divisor $V$ on $X_s$, prove the equivalence of orthogonality to all vertical divisors, $V^2=0$, and membership in the span of the connected components of the scheme-theoretic fibre.}
\pf Write $X_s=\sum d_i\Gamma_i$, $V=\sum a_i\Gamma_i$, and $b_i=a_i/d_i$. Since $X_s\cdot\Gamma_i=0$, expansion gives
\[ V^2=-\sum_{i<j}d_id_j(\Gamma_i\cdot\Gamma_j)(b_i-b_j)^2. \]
All coefficients are nonnegative. Equality holds precisely when $b_i$ is constant on each connected component of the incidence graph. Thus $V$ is a rational combination of the corresponding full fibre components. Each such component is orthogonal to every vertical prime, proving all three implications. \hfill$\square$

\exercise{1.7}{For connected $X_s$, study the intersection form modulo $\Q X_s$, the functional induced by a divisor of generic degree zero, and its vertical correction.}
\pf The formula in 1.6 shows that the induced form on $W=\operatorname{Div}_s(X)\otimes\Q/\Q X_s$ is negative definite. If $\deg D_\eta=0$, the horizontal intersection formula gives $D\cdot X_s=0$, so $w\mapsto D\cdot w$ descends to $W$. Its representing vector $e\in W$ exists uniquely by nondegeneracy. Lift $e$ to a rational vertical divisor $E_0$ and clear denominators: for some $n\ge1$, $E=nE_0$ is integral and $(nD-E)\cdot\Gamma=0$ for every component of $X_s$. Its class is unique modulo $\Q X_s$. \hfill$\square$

\exercise{1.8}{Show that the gcd $d$ of the fibre multiplicities divides $p_a(X_\eta)-1$.}
\pf The integral divisor $F=X_s/d$ has $F^2=0$. Adjunction, including its parity assertion for vertical Cartier divisors, gives $K\cdot F=2p_a(F)-2\in2\Z$. On the other hand $dK\cdot F=K\cdot X_s=2p_a(X_\eta)-2$. Dividing by two proves the assertion. Here $p_a(F)$ is the virtual adjunction genus when needed; because $F\le X_s$, it is also the usual arithmetic genus. \hfill$\square$

\exercise{1.9}{Relate the number $r_\Gamma$ of geometric components and their generic length $e_\Gamma$ to degrees, the canonical degree, and the index of the generic fibre. Include the purely inseparable length calculation.}
\pf (a) Let $k\subset k'\subset k''$ be finite purely inseparable extensions and $F/k$ a field. The tensor products are local Artinian rings. If their residue fields are $F'$ and $F''$, dimension over $F$ gives
\[ \ell(F\otimes_k k')=\frac{[k':k]}{[F':F]},\qquad
\ell(F\otimes_k k'')=\ell(F\otimes_k k')\ell(F'\otimes_{k'}k''). \]
The second formula follows by substitution and the tower law. Infinite algebraic extensions are handled by a finite subextension containing the finitely many coefficients involved.

(b) Pass first to a finite separable splitting field of the geometric components. The $r_\Gamma$ components are conjugate and have equal degrees for a pulled-back invertible sheaf. Next make a finite purely inseparable extension over which their reductions are geometrically reduced. The degree formula for a curve, as the sum of generic lengths times degrees on its reduced components, and (a) show
\[ r_\Gamma e_\Gamma\mid\deg_k L \quad(L\in\operatorname{Pic}(\Gamma)). \]
Indeed after extension the total degree is unchanged, there are $r_\Gamma$ equal contributions, and each has the common generic length $e_\Gamma$. The residue-field and generic-length factors in successive purely inseparable extensions multiply as in (a).

(c) Since $K\cdot X_s=\sum_i d_i\deg_k(\omega_{X/S}|\Gamma_i)$, the integer $d''=\gcd_i(d_ir_i e_i)$ divides $2p_a(X_\eta)-2$. (d) For a closed generic point $P$, its horizontal closure $H$ gives $[K(P):K]=H\cdot X_s=\sum_i d_i\deg_k\mathcal O(H)|\Gamma_i$. Thus $d''$ divides every closed-point degree, hence the index. \hfill$\square$

\exercise{1.10}{Compute the arithmetic genus of a plane curve of degree $n$.}
\pf In $\mathbf P^2$, $K=-3H$, $C=nH$ and $H^2=1$. Adjunction gives $2p_a(C)-2=C(C+K)=n(n-3)$, hence $p_a(C)=(n-1)(n-2)/2$. This computation allows singular and nonreduced effective plane curves. \hfill$\square$

\exercise{1.11}{Prove that $D^2+K\cdot D$ is even for every integral vertical divisor, including divisors with negative coefficients.}
\pf For a vertical prime this is adjunction. Put $q(D)=D^2+K\cdot D$. The identities $q(D+E)=q(D)+q(E)+2D\cdot E$ and $q(-D)=q(D)-2K\cdot D$ propagate evenness from prime divisors to every integral linear combination. This avoids applying effective-divisor adjunction directly to a negative divisor. \hfill$\square$

\exercise{1.12}{If $L$ is generically trivial and has degree zero on each fibre, represent it vertically and show that a bounded power descends to $S$; examine the additional condition from a section.}
\pf A generic trivialization extends rationally and gives $L\simeq\mathcal O(V)$ for a vertical Cartier divisor $V$. \textbf{The printed total-degree condition does not imply descent.} Blow up $\mathbf P^1_R$ at a rational point of its special fibre, with exceptional curve $E$. Then $L=\mathcal O(E)$ is generically trivial and $E\cdot X_s=0$, but $E^2=-1$, so $L$ is not pulled back from $\Spec R$. The gcd of the fibre multiplicities is one, so even the prescribed power does not repair this example.

The intended condition is $\deg L|\Gamma=0$ for \emph{each irreducible fibre component}. The generic smooth geometrically connected fibre implies connected fibres. By 1.6, $V_s=a_sX_s/d_s$ with $a_s\in\Z$ and $d_s$ the gcd of its multiplicities. Smoothness over a dense open gives $d_s=1$ except at finitely many $s$, so $d=\operatorname{lcm}_s d_s$ is finite. Then $dV=\pi^*B$ for a divisor $B$ on $S$ and $L^d=\pi^*\mathcal O_S(B)$. If a section $\sigma$ has $\sigma^*L^d$ trivial, then $\mathcal O_S(B)$ is trivial, proving (c) under this corrected hypothesis. \hfill$\square$

\exercise{1.13}{For a cohomologically flat arithmetic surface with $\pi_*\mathcal O_X=\mathcal O_S$, descend fibrewise trivial line bundles and compare two line bundles using a section.}
\pf Cohomological flatness gives $H^0(X_s,\mathcal O_{X_s})=k(s)$. If $L|X_s$ is trivial for every $s$, cohomology and base change show that $M=\pi_*L$ is invertible, with $M\otimes k(s)=H^0(X_s,L_s)$. The evaluation map $\pi^*M\to L$ is an isomorphism on every fibre; Nakayama makes it surjective on $X$, and a surjection of invertible sheaves is an isomorphism. For (b), apply this to $L=L_1\otimes L_2^{-1}$. The section gives $M=\sigma^*L$, so the stated equality on the section is exactly the condition $M\simeq\mathcal O_S$. Necessity follows by restriction. \hfill$\square$

\exercise{1.14}{Compare the reduced divisor $Z_0=\sum_{i\le m}\Gamma_i$ and $Z=\sum_{i\le m}(d_i/d)\Gamma_i$ by a chain adding components of positive intersection, and deduce injectivity on global functions.}
\pf \textbf{For an arbitrary partial fibre both claims are false.} Start with $\mathbf P^1_R$, blow up a rational special point, and then blow up the intersection of the two resulting fibre components. The new exceptional curve $B$ has multiplicity two and $B^2=-1$, while the other two components have multiplicity one. Take $m=1$, $\Gamma_1=B$, so $d=1$, $Z_0=B$, $Z=2B$. The only possible step adds $B$, but $B^2<0$. Moreover
\[ 0\longrightarrow\mathcal O_B(-B)\longrightarrow\mathcal O_{2B}
\longrightarrow\mathcal O_B\longrightarrow0 \]
has $H^0(\mathcal O_B(-B))=H^0(\mathbf P^1_k,\mathcal O(1))\ne0$, contradicting (b).

For the \emph{full} primitive fibre $F=X_s/d$, the assertions are true, allowing an empty chain if $F$ is reduced. Here is the required algorithm. Start at $A=(X_s)_{\mathrm{red}}\le F$. If $A<F$ and every $A\cdot\Gamma_i\le0$, then $A\cdot X_s=0$ forces all these intersections to vanish. By 1.6, $A=cF$; integrality and primitivity force $c$ to be an integer, impossible for $0<A<F$. Thus some $A\cdot\Gamma>0$. Its coefficient is smaller than that in $F$, since otherwise $(F-A)\cdot\Gamma\ge0$ would contradict $F\cdot\Gamma=0$. Add $\Gamma$ and repeat; bounded coefficients ensure termination at $F$. Each step has kernel $\mathcal O_\Gamma(-A)$ of negative degree, hence no global sections. Composing proves $H^0(F,\mathcal O_F)\hookrightarrow H^0(F_{\mathrm{red}},\mathcal O)$. For a proper negative-definite configuration the analogous valid endpoint is its fundamental cycle, rather than the restriction of $F$. \hfill$\square$

\exercise{1.15}{Compute canonical degrees and self-intersection of the graph of a finite map of smooth proper curves; compute its intersections with the diagonal and specialize to Frobenius.}
\pf We use smooth geometrically connected curves, as required by the displayed genus formulas. For curves with larger constant fields replace $2g-2$ by their actual canonical degrees. The product differential sequence gives
\[ \omega_{C\times E/k}=p_1^*\omega_C\otimes p_2^*\omega_E,
\qquad \omega_{C\times E/k}|\Gamma_f=\omega_C\otimes f^*\omega_E. \]
Its degree is $2g(C)-2+\deg(f)(2g(E)-2)$. Since $\Gamma_f\simeq C$, adjunction gives $\Gamma_f^2=\deg(f)(2-2g(E))$ and in particular $\Delta^2=2-2g(C)$.

For $C=E$, a rational fixed point $y$ and uniformizer $t$ give the local quotient $\mathcal O_{C,y}/(f^*t-t)$, so its length is the local intersection number if $f\ne\operatorname{Id}$. The print's word ``automorphism'' should be ``local homomorphism'' when $f$ is not invertible. For $q$-Frobenius, its geometric fixed points are exactly $C(\mathbf F_q)$; its tangent map is zero. Thus $f^*t-t$ has a nonzero linear term at every fixed point, every intersection has length one, and $\Gamma_f\cdot\Delta=\#C(\mathbf F_q)$. \hfill$\square$

\section*{9.2 Intersection and morphisms}

\exercise{2.1}{Factor a proper birational morphism between regular two-dimensional locally Noetherian schemes into point blow-ups, and prove projectivity over a quasi-projective target.}
\pf For a Noetherian target this is the proper version of the surface factorization theorem. The proof of Lemma 2.1 is local on the target: at a point with a positive-dimensional fibre, the pullback of its maximal ideal is invertible, so the morphism factors through its blow-up. The resulting morphism has one fewer exceptional prime. Properness and quasi-compactness make the exceptional set finite; iteration ends with a proper quasi-finite birational morphism to a normal scheme, hence an isomorphism. This argument does not require the original morphism to have been given as projective. Each point blow-up is projective over a target quasi-projective over an affine Noetherian scheme: twist its coherent ideal by an ample invertible sheaf to make it globally generated, without changing the blow-up. Composition proves (b).

For a merely locally Noetherian target, (a) is interpreted locally, or with a locally finite sequence of centers. A globally finite sequence is false without quasi-compactness: on the disjoint union of countably many copies of $\mathbf A^2$, blow up one origin in each copy. The resulting morphism is proper and birational, but no finite list of individual closed-point blow-ups produces it. The hypothesis in (b) restores quasi-compactness. \hfill$\square$

\exercise{2.2}{Describe a multiplicity-one exceptional divisor $E$ in a special fibre.}
\pf Set $k'=H^0(E,\mathcal O_E)$ and $c=[k':k(s)]$. Then $E\simeq\mathbf P^1_{k'}$ and $E^2=-c$. Writing $X_s=E+\sum d_i\Gamma_i$, we obtain
\[ c=\sum d_i(E\cdot\Gamma_i). \]
Every nonzero summand $E\cdot\Gamma_i$ is a positive multiple of $c$, because it is the degree over $k(s)$ of a divisor on $E$. Hence exactly one summand occurs, its multiplicity is one, and its intersection degree is $c$. The local length and residue degree are consequently both one over $k'$. Thus $E$ meets exactly that component, transversally at one $k'$-rational point. \hfill$\square$

\exercise{2.3}{Reprove the exceptional-divisor formula of Proposition 2.5 using Proposition 2.23.}
\pf At the blown-up point choose regular parameters $u,v$. The exceptional curve is $\operatorname{Proj}\operatorname{gr}_{\fm}A=\mathbf P^1_{k(x)}$. The strict transform of $u=0$ meets it at a single $k(x)$-rational point with length one, as the chart $u=vw$ shows. Proposition 2.23 gives $\pi^*\operatorname{div}(u)=\widetilde{\operatorname{div}(u)}+E$ locally. Restriction of its associated line bundle to $E$ is trivial, whereas the strict transform restricts to $\mathcal O_E(1)$. Therefore $\mathcal O(E)|E=\mathcal O_E(-1)$ and $E^2=-[k(x):k(s)]$. Also $H^0(E,\mathcal O_E)=k(x)$. \hfill$\square$

\exercise{2.4}{Compute the Picard group and canonical discrepancy of a birational morphism $f:X\to Y$ of regular fibred surfaces with exceptional primes $E_1,\ldots,E_n$.}
\pf For a single point blow-up, restriction off $E$ and extension of Cartier divisors on a regular surface give
\[ \operatorname{Pic}(\widetilde Y)=\pi^*\operatorname{Pic}(Y)\oplus\Z[E]. \]
The sum is direct since $E^2\ne0$. Moreover $\omega_{\widetilde Y/S}=\pi^*\omega_{Y/S}\otimes\mathcal O(E)$ by Proposition 2.24. Factor $f$ into point blow-ups. Successive total transforms of the new exceptional curves form a basis of the relative Picard group. Their change to the final strict transforms is triangular with diagonal entries one, proving (a) with these $n$ primes as basis. At each step the discrepancy is the pullback of the preceding effective discrepancy plus the new exceptional curve. Old coefficients stay positive, and each new coefficient is one plus a sum of nonnegative old contributions. Thus $\omega_{X/S}=f^*\omega_{Y/S}\otimes\mathcal O(\sum a_iE_i)$ with every $a_i>0$. \hfill$\square$

\exercise{2.5}{Prove $\mu_x(D+E)=\mu_x(D)+\mu_x(E)$ on a regular surface.}
\pf The associated graded ring of its regular local ring is a polynomial ring over $k(x)$, hence a domain. The product of the nonzero initial forms of equations $f,g$ has degree $\operatorname{ord}(f)+\operatorname{ord}(g)$ and is nonzero. Consequently $\operatorname{ord}(fg)$ is exactly this sum. \hfill$\square$

\exercise{2.6}{Separate two sections by blow-ups, counting the exceptional curves, and examine normal crossings of their union with the affected fibres.}
\pf A section passes through a smooth point of the fibre: the base uniformizer is not in $\fm_x^2$, since its image under the retraction to the local DVR of the section has order one. Near an intersection of the sections write the completed local equations as $z=0$ and $z=a(t)$, with $\operatorname{ord}_t a=r=i_x(D,F)$. Blow up their common point. Their strict transforms either separate if $r=1$, or meet with intersection multiplicity $r-1$. Repetition uses exactly $r$ point blow-ups. The charts $z=tz_1$ show that the last intersections with the exceptional fibre are distinct and transverse, and the resulting total divisor has normal crossings near all the centers. The same calculation is valid before completion by faithful flatness. The centers over different intersection points are disjoint, so the count is $\sum_x i_x(D,F)$.

\textbf{The unrestricted global normal-crossings claim is false.} Over $R=k[[t]]$, $\operatorname{char}k=0$, the regular cubic surface $Y^2Z=X^3+tZ^3$ has a cuspidal special fibre. The two sections with affine coordinates
\[ (1,\sqrt{1+t}),\qquad(1+t,\sqrt{(1+t)^3+t}) \]
meet with multiplicity one at the smooth point $(1,1)$. A single point blow-up resolving their triple intersection with the fibre leaves the cusp unchanged. A blow-up at the cusp instead leaves that triple intersection unchanged. Thus the asserted count one cannot yield global normal crossings. The counted construction proves the statement if the selected fibres already have normal crossings away from the section intersections; in general additional centers resolving the other singularities are necessary. \hfill$\square$

\exercise{2.7}{Produce the local hypersurface presentation $\mathcal O_{X,x}\simeq\mathcal O_{Z,z}/(\prod u_i^{d_i}-ta)$ with $Z_s$ regular of dimension two, and preserve the multiplicities $\mu_i=\mu_x(\Gamma_i)$.}
\pf Put $A=\mathcal O_{X,x}$ and write $t=b\prod f_i^{d_i}$ in this factorial regular local ring. If the fibre is regular at $x$, only one component occurs, with multiplicity one. Take locally $Z=X\times\mathbf A^1$, coordinate $v$, and lift $f_1$ as $u_1=f_1+v$; with a lifted unit $a=b^{-1}$ the required equation is $v=0$ up to a unit. The orders are one.

Otherwise $t\in\fm_A^2$ and $A/tA$ has embedding dimension two. Embed a neighborhood in affine space over $S$. In the regular local ring of its special fibre choose generators of the immersion ideal whose independent linear terms cut its embedding dimension down to two. There are $N-2$ such equations in a relative $N$-dimensional ambient space. Their common zero scheme $Z$ is flat over $S$ by the regular-sequence criterion, and $Z_s$ is regular of dimension two at $z$. Thus $B=\mathcal O_{Z,z}$ is regular of dimension three, $t$ is a parameter of $B$, and $A=B/(h)$ with $h$ a parameter. Since $t$ has zero linear term in $A$, the linear term of $h$ is a nonzero multiple of that of $t$.

Lift the $f_i$ to $u_i\in B$ and $b^{-1}$ to a unit $a$. The element $F=\prod u_i^{d_i}-ta$ belongs to $(h)$. Its linear term is $-\overline a\,t$: the product has order at least two because $\sum d_i\mu_i=\operatorname{ord}_A(t)\ge2$. Hence $F/h$ is a unit and $A=B/(F)$. The associated graded map identifies $\operatorname{gr}_{\fm_A}A$ with $\operatorname{gr}_{\fm_B}B/(\operatorname{in}(t))$. Choose each lift with the same first nonzero homogeneous term, involving the other two parameters. Its order in $B/tB$ is then exactly $\mu_i$, proving (b). Shrinking the affine schemes realizes these local constructions by finite-type schemes as requested. \hfill$\square$

\exercise{2.8}{Split fibre components by finite etale residue extensions, compute intersections after splitting, and test product divisibility by the geometric invariants.}
\pf (a) For a finite Galois splitting field $k'/k$, the $r_\Gamma$ components of $\Gamma_{k'}$ are permuted transitively by Galois and are isomorphic as $k$-schemes. (b) Lift the separable minimal polynomial of a primitive generator of $k'/k$ to a monic polynomial over $R$. Its quotient $R'$ is finite free and local, since its special fibre is the field $k'$; the derivative is a unit, so it is finite etale. Being connected and normal it is a DVR. For a chosen component $\Gamma_i$,
\[ p_*\Gamma_i=\frac{[k':k]}{r_\Gamma}\,\Gamma. \]

(c) Pullback is $p^*\Delta=\sum_j\Delta_j$, and the projection formula, with the residue-field conversion of degrees, gives
\[ \Gamma\cdot\Delta=r_\Gamma\sum_j\Gamma_i\cdot\Delta_j. \]
The sum is a degree on $\Gamma_i$, whose geometric generic multiplicity remains $e_\Gamma$; hence $r_\Gamma e_\Gamma$ divides $\Gamma\cdot\Delta$. Each term is also a degree on $\Delta_j$, whose generic multiplicity is $e_\Delta$; hence $r_\Gamma e_\Delta$ divides it. Interchanging the two curves gives the other two divisibilities $r_\Delta e_\Gamma$ and $r_\Delta e_\Delta$.

(d) Products of the two $r$'s or the two $e$'s need not divide. Let $R=k[[t]]$, take a degree-two separable extension $l/k$, and blow up a rational special point of $\mathbf P^1_{l[[t]]}$, viewed over $R$. The two fibre components have $r=2$ and intersect with degree two, which is not divisible by four. For the $e$ example use instead a purely inseparable extension $l/k$ of degree $p$. Both components then have $e=p$ and intersect with degree $p$, not $p^2$. These total surfaces are regular and projective flat over $R$; geometrical connectedness is not required in the question. \hfill$\square$

\exercise{2.9}{Study the three fibre gcd invariants under a point blow-up and compute the changed intersection matrix.}
\pf Put $\mu_i=\mu_x(\Gamma_i)$ and $c=[k(x):k]$. The exceptional $E=\mathbf P^1_{k(x)}$ has $r_E$ equal to the separable degree of $k(x)/k$ and $e_E$ its inseparable degree, so $r_Ee_E=c$. Proposition 2.23 gives
\[ d_E=\sum_i d_i\mu_i,\qquad \pi^*\Gamma_i=\widetilde\Gamma_i+\mu_i E. \]
Strict transforms have the same geometric component invariants, since their function fields are unchanged. By 2.8 applied to $\widetilde\Gamma_i,E$, and by the formula below, $r_ie_i\mid\mu_i r_Ee_E$.

For the stronger separable assertion $r_i\mid\mu_i r_E$, extend to a separable closure. There are $r_E$ lifts of $x$ and $r_i$ conjugate components above $\Gamma_i$. Sum the orders of the local equations over the lifts of $x$. The total is $r_E\mu_i$, and transitivity makes the contribution from each of the $r_i$ components equal. This proves the divisibility, including when these integers have residue-characteristic factors.

Let $d=\gcd d_i$, $d'=\gcd(d_ir_i)$ and $d''=\gcd(d_ir_ie_i)$. The preceding formulas show that each old gcd divides the corresponding new term $d_E$, $d_Er_E$, or $d_Er_Ee_E$. Since all old terms remain, the three gcd's are unchanged. A common regular refinement and factorization into point blow-ups give birational invariance.

Finally $\pi^*\Gamma_i\cdot E=0$ and $E^2=-c$ imply
\[ \widetilde\Gamma_i\cdot E=\mu_i c,\qquad
\widetilde\Gamma_i\cdot\widetilde\Gamma_j=\Gamma_i\cdot\Gamma_j-\mu_i\mu_jc. \]
The second formula includes $i=j$. \hfill$\square$

\exercise{2.10}{Bound the degree of a closed generic point by the multiplicities at its specializations, and compute the fibre multiplicity at a closed point.}
\pf Its horizontal closure $H$ is finite flat over the local DVR: it is proper, has zero-dimensional fibres, and is torsion-free. Thus
\[ [K(P):K]=H\cdot X_s=\sum_{x\in H\cap X_s}i_x(H,X_s)[k(x):k(s)]. \]
In a regular two-dimensional local ring, intersection multiplicity is at least the product of orders. One can see this by blowing up once: the total-transform formula and projection give $i_x(H,X_s)\ge\mu_x(H)\mu_x(X_s)$. Since $\mu_x(H)\ge1$, this proves the bound. Taking orders in $t=b\prod f_i^{d_i}$, using 2.5, gives $\mu_x(X_s)=\sum_i d_i\mu_x(\Gamma_i)$. \hfill$\square$

\exercise{2.11}{Over a Henselian DVR, construct a closed generic point of the predicted degree through $x$ when the residue field is sufficiently large; examine the displayed $p$-adic example.}
\pf (a) The finite integral scheme $H=\overline{\{P\}}$ is connected and finite over the Henselian local base. A finite algebra over that base is the product of its local factors. Since $H$ is integral it has one factor, hence only one closed special point.

(b) In the two-dimensional regular local ring $\mathcal O_{Z_s,z}$ of 2.7, each nonzero degree-$\mu_i$ initial form has at most $\mu_i$ linear factors over $k(x)$. There are $\#k(x)+1$ directions in $\mathbf P^1(k(x))$ when the field is finite. If $\#k(x)\ge\sum_i\mu_i$, choose a direction avoiding all these factors; the same finite avoidance works over an infinite field. Lift its linear equation to $u$. This proves the requested nondivisibility, and also the stronger transversality to every tangent cone needed in (c).

(c) The image of $u$ is a parameter in the regular ring $\mathcal O_{X,x}$. In the singular-fibre case this follows from the cotangent identification in 2.7; in the smooth-fibre case it follows from the independent chosen direction. Thus its local zero divisor is a single regular branch $D$. In $\mathcal O_{Z_s,z}$ the chosen line has intersection $\mu_i$ with $u_i=0$, so
\[ i_x(D,X_s)=\sum_i d_i\mu_i=\mu_x(X_s). \]
This is finite, so the branch is horizontal. Let $P=D_\eta$. Part (a) shows that $x$ is its only specialization, and the horizontal degree formula gives $[K(P):K]=\mu_x(X_s)[k(x):k(s)]$. Mere nondivisibility by the full $u_i$ would not suffice; the tangent-cone avoidance above is the required choice.

(d) The displayed hypersurface is flat and projective. Its generic fibre is smooth: on $w\ne0$ the derivative in $w$ is nonzero in characteristic zero, and at infinity the binary form has distinct roots. The special fibre consists of $p+1$ distinct lines through $x$. Away from $x$ it is smooth; at $x$ the defining equation has linear term $p$, so the total surface is regular there as well. Its multiplicity is $p+1$.

\textbf{The printed bound $2p$ is false for $p\ge3$.} Blow up $x$. The exceptional curve has multiplicity $p+1$; the $p+1$ strict lines, of multiplicity one, meet it at all its $\mathbf F_p$-rational points. If the strict horizontal closure meets a rational point of $E$, 2.10 gives degree at least $p+2$. If it meets a nonrational point, it gives at least $2(p+1)$. Hence the valid lower bound is $p+2$, still strictly greater than $\mu_x(X_s)$.

It is sharp. In the affine chart $w=1$, impose $v=u+u^2$. The resulting polynomial is
\[ G(u)=p+u^{p+1}(1+u)\bigl(1-(1+u)^{p-1}\bigr). \]
Modulo $p$ its first nonzero term is $-(p-1)u^{p+2}$. Weierstrass preparation in $\Z_p[[u]]$ factors $G$ into a unit and a distinguished polynomial of degree $p+2$, whose constant coefficient has valuation one. This polynomial is Eisenstein. A root $\alpha$ gives $[\Q_p(\alpha):\Q_p]=p+2$ and a point $(\alpha,\alpha+\alpha^2,1)$ specializing to $x$. For $p>2$, $p+2<2p$. For $p=2$, the binary form still has three distinct lines and the valid bound four equals $2p$. \hfill$\square$

\exercise{2.12}{Compute canonical intersections and the genus change for the strict transform of a vertical prime under a point blow-up.}
\pf With $\mu=\mu_x(\Gamma)$, $\pi^*\Gamma=\widetilde\Gamma+\mu E$, $K_{\widetilde X/S}=\pi^*K_{X/S}+E$, and $\pi^*D\cdot E=0$, expansion gives
\[ K_{\widetilde X/S}\cdot\widetilde\Gamma=K_{X/S}\cdot\Gamma-\mu E^2,
\qquad \widetilde\Gamma^2=\Gamma^2+\mu^2E^2. \]
Add the equations and apply adjunction, using $E^2=-[k(x):k(s)]$, to obtain
\[ p_a(\widetilde\Gamma)=p_a(\Gamma)-\frac{\mu^2-\mu}{2}[k(x):k(s)]. \]
\hfill$\square$

\exercise{2.13}{Prove Abhyankar's theorem: an exceptional prime over a regular excellent scheme is birationally a projective-space bundle over an integral variety.}
\pf Let $\nu$ be the DVR valuation at its generic point $\xi$, and $y_0$ its center. Every nonzero ideal pulled back to this DVR is principal. The universal property of a blow-up therefore lifts the map at $\xi$ through the blow-up of $\overline{\{y_0\}}$. Repeat with its successive centers $y_i$. The valuation-ring approximation argument of Exercise 8.3.14 says that the local rings of these successive quadratic transforms reach $\mathcal O_{X,\xi}$ after finitely many steps: a divisorial valuation has finite transcendence-degree defect zero, and its finitely many residue generators and a uniformizer occur in a finite stage. Thus there is a last center of local dimension at least two; call it $y_r$ (with the initial center included in the indexing).

All these finitely many center closures are generically regular. Excellence makes their regular loci open; remove the images of their bad closed subsets, after shrinking around the original generic center. Inductively the ambient schemes are regular as well, since blowing up a regular center in a regular scheme is regular. At the next stage the center $y_{r+1}$ is a codimension-one point with local DVR equal to $\mathcal O_{X,\xi}$. The birational map is consequently an isomorphism on neighborhoods of these generic points: in affine presentations the finitely many generators and denominators of both inverse maps are units there, and can be inverted. This is the sense of the local isomorphism required in (d).

If the last center closure is $F_r$ of codimension $c\ge2$, the last exceptional divisor is $\mathbf P(N_{F_r/Y_r})$, generically a $\mathbf P^{c-1}$-bundle. Its generic point is the center of $\nu$, so it has the same function field as $E$. A vector bundle is trivial on a dense open of $F_r$; hence $E$ is birational to $\mathbf P^{c-1}_{F_r}$, with $d=c-1\ge1$. \hfill$\square$

\section*{9.3 Minimal surfaces}

\exercise{3.1}{Characterize relatively minimal conic models, elementary transformations at split fibres, and existence of an absolute minimal regular model.}
\pf The generic conic has arithmetic genus zero. By 1.8, the gcd of each connected fibre's multiplicities is one. On a reducible connected fibre every component has negative self-intersection; since $K\cdot X_s=-2$, some component also has negative canonical intersection. Proposition 3.10 makes it exceptional. Conversely an irreducible fibre has self-intersection zero and contains no exceptional curve. Its single multiplicity is one by the gcd assertion, so it is integral. This proves (a).

If $X_s=\mathbf P^1_k$, blow up a rational point. The strict fibre and new exceptional curve are both $\mathbf P^1_k$ with square $-1$. Contract the former. This gives a different regular model of the same generic conic which is not dominated by $X$: the induced birational map has a base point at the blown-up point. Thus $X$ is not minimal. For the other assertion in (b), a point over $\operatorname{Frac}\widehat{\mathcal O}_{S,s}$ specializes on the proper model to a $k$-rational point. Its section passes through the smooth locus. An integral arithmetic-genus-zero fibre with a smooth rational point is $\mathbf P^1_k$ (the genus-zero curve criterion). Conversely a rational smooth point lifts over the complete DVR by formal smoothness and properness.

We recall explicitly the elementary-transformation argument used in (c). In a factorization of a birational map between relatively minimal regular conic models, cancel consecutive inverse blow-ups and commute blow-ups at disjoint points. The first contraction which is not a cancellation must contract the strict transform of an original fibre component. The original fibre is integral and has square zero. If points of degrees $c_1,\ldots,c_l$ on its successive strict transforms have been blown up, its square is $-\sum c_i$. To be exceptional it must be a rational curve over its field of constants and have square minus that constant-field degree. An integral genus-zero fibre has constants $k(s)$; hence $\sum c_i=1$. There was exactly one blow-up, at a rational point, and the strict fibre is $\mathbf P^1_{k(s)}$. Thus this first change of model is precisely the elementary transformation above. This also justifies the cancellation argument without assuming all birational maps are morphisms.

Consequently a relatively minimal but nonminimal model has a common intermediate model with fibre $E_1+E_2$ consisting of these two exceptional curves. Adjunction gives
\[ -2=K_Z\cdot E_1+K_Z\cdot E_2
=-[H^0(E_1,\mathcal O):k]-[H^0(E_2,\mathcal O):k], \]
so both constant fields are $k$ and both curves are $\mathbf P^1_k$. Part (b) gives a point over the completed fraction field. Finally relatively minimal models exist by successive exceptional contractions. If every completed conic is anisotropic, the preceding argument allows no elementary transformation, so such a model is minimal. If one completed conic has a point, every relatively minimal model has a split fibre there and fails minimality. This proves (d). \hfill$\square$

\exercise{3.2}{Prove minimality of a relatively minimal projective surface with a nonzero pluricanonical section; explain the rational and ruled exceptions.}
\pf If it were not minimal, the common-resolution and cancellation argument in Theorem 3.21 produces a regular $Z$ with two distinct intersecting exceptional curves $E_1,E_2$. Pullback of a nonzero $q$-canonical section and the effective discrepancies of 2.4 give an effective divisor
\[ qK_Z=a_1E_1+a_2E_2+D,\qquad a_i\ge0,\quad D\ge0, \]
where $D$ contains neither curve. Put $c_i=-E_i^2=[H^0(E_i,\mathcal O):k]>0$ and $I=E_1\cdot E_2>0$. Since $K_Z\cdot E_i=E_i^2$, intersection gives
\[ (a_1-q)c_1\ge a_2 I,\qquad(a_2-q)c_2\ge a_1 I. \]
The positive number $I$ is divisible by each $c_i$, being a degree on $\mathbf P^1_{k_i}$. Thus $a_1-q\ge a_2$ and $a_2-q\ge a_1$, a contradiction. This proves (a)--(c).

Here ``minimal'' means dominated by every regular birational model, stronger than having no exceptional curves. For a rational surface, such a universal model would be dominated by $\mathbf P^2$ and by $\mathbf P^1\times\mathbf P^1$. Neither has exceptional curves, so the dominating morphisms would both be isomorphisms, impossible since their Picard ranks differ. A ruled surface admits elementary transformations of its $\mathbf P^1$ fibres as in 3.1; relatively minimal representatives related by a nontrivial such transformation cannot dominate one another as models. Thus these birational classes have no universal minimal representative, though they do have relatively minimal ones. \hfill$\square$

\exercise{3.3}{Give a relatively minimal conic model whose finite etale base change is not relatively minimal.}
\pf Let $R=k[[t]]$, $\operatorname{char}k\ne2$, and $a\in k^*$ nonsquare. Set
\[ X=\operatorname{Proj}R[u,v,w]/(u^2-av^2+tw^2). \]
The generic fibre is a smooth conic. The special fibre is integral over $k$, so 3.1 gives relative minimality. At its only geometrically singular point $[0:0:1]$ the defining equation has nonzero linear term $t$ in the regular ambient local ring; the total surface is regular. After the finite etale extension $R'=k(\sqrt a)[[t]]$, the special fibre becomes two lines meeting transversally, each of multiplicity one. Their intersections have degree one and each has square $-1$ since it intersects the whole fibre in zero. They are exceptional, so the base-changed model is not relatively minimal. \hfill$\square$

\exercise{3.4}{Relate automorphisms of a normal surface and of its minimal desingularization.}
\pf If $f:X\to Y$ is the minimal desingularization, an automorphism $\alpha$ of $Y$ makes $\alpha\circ f$ another minimal desingularization. Uniqueness gives a unique lift $\widetilde\alpha$ to $X$; uniqueness also makes lifting compatible with composition. Conversely let an automorphism of $X$ preserve the exceptional locus setwise. Its connected components are exactly the nontrivial fibres of $f$ over the singular points, since $f_*\mathcal O_X=\mathcal O_Y$. Thus $f\circ\beta$ is constant on every fibre of $f$ and descends, by the normal contraction descent criterion, to a morphism $\alpha:Y\to Y$. Apply the same argument to $\beta^{-1}$ to obtain its inverse. The two constructions are inverse, and equality on the dense regular open proves uniqueness. \hfill$\square$

\exercise{3.5}{Seek a connected finite etale Galois refinement of an etale surjection $S'\to S$.}
\pf \textbf{The assertion requires $S'\to S$ finite (and a connected component dominating $S$).} For a counterexample even with integral $S'$, take $k=\mathbf F_p$, $S=\mathbf A^1_k$ and the Artin--Schreier cover $t=u^p-u$. Remove the point $u=0$ from the covering affine line. The resulting integral $S'\to S$ remains etale and surjective: above $t=0$ the other $p-1$ points remain. It is not finite. If a nonempty finite etale $T\to S$ factored through $S'$, its map to $\mathbf A^1_u$ would be proper (the target is separated over $S$) and have dense image. It would therefore be surjective, contradicting the omitted point.

For a connected finite etale $S'\to S$, choose a finite Galois extension $N/K(S)$ containing its function field. Normalize $S$ in $N$. This is finite etale: equivalently use the finite etale fundamental-group action on a geometric generic fibre and take the kernel of its permutation representation. The connected Galois cover corresponding to that kernel dominates $S'$. In the exercise's product construction all conjugates are finite etale; the component containing the chosen generic embedding has function field the Galois closure, is finite etale, and maps surjectively onto $S'$. Its automorphism group is the generic Galois group because the base is normal. These observations prove (a)--(c) in the finite case. For a general etale surjection one can use etale local trivializations, but one cannot assert a finite refinement over the original base or even over a Zariski neighborhood of every point. Exercise 3.6 below therefore supplements the finite-cover argument with a numerical proof for arbitrary etale maps. \hfill$\square$

\exercise{3.6}{Give an etale-descent proof that minimal arithmetic surfaces of generic arithmetic genus at least one remain minimal under etale base change.}
\pf First use a connected finite Galois cover $T\to S$ with group $G$. Uniqueness of the minimal model of $X_T$ extends its generic $G$-action to its minimal model $Z$. The contraction $X_T\to Z$ is equivariant, again by uniqueness. Projective Galois descent is effective: tensor the translates of an ample line bundle to obtain an invariant polarization and descend its section algebra. It produces $Y=Z/G$ over $S$, with
\[ Z\simeq Y\times_S T. \]
In particular $Y$ is regular, since regularity descends under etale faithfully flat maps, and $X\to Y$ is a birational contraction. Minimality of $X$ forces this morphism to be an isomorphism. Therefore $X_T=Z$ is minimal, and minimality descends to every intermediate etale cover. This establishes the requested quotient and normality statements as well: the displayed base change is the regular scheme $Z$, not merely an unexplained consequence of normality.

For an arbitrary etale $S'\to S$, the finite global refinement invoked in the print need not exist, as 3.5 shows. Use instead the following numerical local check, which does not require that refinement. An exceptional curve on $X_{S'}$ maps finitely onto a fibre component $C$ of $X$. Etale base change preserves the canonical sheaf; its negative canonical degree would imply $K_X\cdot C<0$ by the degree projection formula. But $K_X$ is nef by Corollary 3.26. This contradiction proves minimality on every component of $S'$. \hfill$\square$

\exercise{3.7}{Determine the minimal desingularization of $u^2w+v^3+t^6w^3=0$ over a DVR of residue characteristic different from two and three.}
\pf Let $\mathcal E/R$ be the smooth constant-coefficient elliptic scheme
\[ U^2W+V^3+W^3=0,\qquad o=[1:0:0], \]
and set $Z=\operatorname{Bl}_{o_s}\mathcal E$. The generic change of coordinates
\[ [U:V:W]\longmapsto[t^3U:t^2V:W] \]
identifies its generic fibre with that of $X$. We check extension across the sole base point $o_s$. On $U=1$, put $q=V/U$, $r=W/U$. The equation $r+q^3+r^3=0$ gives $r=q^3h$ with $h$ a unit near $o$. The base ideal is
\[ (t^3,t^2q,q^3h). \]
On the two charts of the blow-up of $(t,q)$, it is generated by $t^3$ when $q=tb$, and by $q^3$ when $t=qa$. After cancelling these factors the homogeneous coordinates are respectively $[1:b:b^3h]$ and $[a^3:a^2:h]$. They have no common zero, so give a morphism $f:Z\to X$.

The special fibre of $Z$ is $\widetilde{\mathcal E}_s+P$, with $P\simeq\mathbf P^1_k$ and both self-intersections $-1$. The genus-one component is contracted to $[0:0:1]$; $P$ maps birationally onto the cuspidal cubic, by the displayed parametrization. There are no other contracted curves. The hypersurface $X$ is regular away from that one point and normal by the hypersurface Serre criterion. Proper birationality and the absence of other contracted curves make $f$ an isomorphism off it. Thus $Z$ is a desingularization. Its only curve exceptional over $X$ has genus one, and so is not an exceptional divisor in Castelnuovo's sense. Proposition 3.32 therefore makes $Z\to X$ the minimal desingularization. Notice that $P$, although a $(-1)$-curve on $Z$, is not contracted by this morphism. \hfill$\square$

\exercise{3.8}{On a minimal arithmetic surface with $p_a(X_\eta)=1$, represent the canonical sheaf by a rational combination of complete fibres.}
\pf First assume the generic constant field is $K(S)$ (in particular, assume geometric connectedness in the smooth case). Its canonical degree is zero and its one-dimensional space of differentials has a nonzero section with no zeros. A rational extension gives $\omega_{X/S}=\mathcal O(V)$ with $V$ vertical. Minimality makes $K$ nef. Since $K\cdot X_s=0$, each $K\cdot\Gamma_i\ge0$ is zero. On each connected fibre, 1.6 yields $V_s\in\Q X_s$, proving the assertion when the fibres are connected. The coefficients of the prime components are integral; no global effectivity is asserted.

\textbf{Arithmetic genus one alone does not impose that constant-field hypothesis.} The always valid statement is a rational combination of the connected components of the fibres, or of the full fibres over the normal Stein base. The canonical sheaf of the generic curve is still trivial over a larger constant field $L$: duality for the finite field extension contributes a one-dimensional $L$-space. But the kernel of the intersection form can have more than one generator above a point of $S$.

Here is a global counterexample to the full-fibre assertion without the extra hypothesis. Over an algebraically closed field of characteristic zero choose a connected etale double cover $q:S_1\to S$ of smooth projective curves. Choose a line bundle $M$ of sufficiently large odd degree on $S_1$. A general Weierstrass fibration with coefficients in $M^4$ and $M^6$ has simple discriminant zeros; its total surface $\pi:X\to S_1$ is regular, its fibres are smooth or irreducible nodal cubics, and $\omega_{X/S_1}=\pi^*M$. (At a simple discriminant zero the local smoothing parameter has order one, so the nodal total space is regular.) View $X$ over $S$. It is minimal, and the generic curve over $K(S)$ has Euler characteristic zero, hence arithmetic genus one. Since $q$ is etale, $\omega_{X/S}=\pi^*M$. Every prime fibre component over $S$ has multiplicity one. A divisor in $\bigoplus_s\Q X_s$ which is integral is therefore pulled back from an integral divisor on $S$. Pullback by the section of $\pi$ would then give $M=q^*B$, impossible because $\deg M$ is odd and $\deg q^*B$ is even. This proves the need to use connected fibre components or the Stein base in the general printed setting. \hfill$\square$

\exercise{3.9}{Compare global functions on $X$ and on an open $U$ containing the generic fibre and meeting every fibre; compare canonical sections in genus one.}
\pf (a) A regular function on $U$ restricts to a constant $a\in K$ on the proper geometrically integral generic curve. Every nonempty open $U\cap X_s$ contains the generic point of some component $\Gamma$. There its valuation is $d_\Gamma\nu_s(a)$, so regularity forces $\nu_s(a)\ge0$. Thus $a\in\mathcal O_S(S)$. Conversely such a function is regular on $X$, proving both rings equal.

(b) \textbf{Minimality is missing.} On the blow-up $b:Z\to\mathcal E$ of a rational point of the special fibre of a smooth elliptic scheme over a DVR, let $E$ be exceptional, $C$ the strict fibre and $\eta$ a nowhere-zero invariant differential on $\mathcal E$. As a section of $\omega_{Z/R}$, $b^*\eta$ has divisor $E$, whereas $\operatorname{div}(t)=C+E$. Hence $t^{-1}b^*\eta$ is regular on $U=Z\setminus C$, which contains the whole generic fibre and meets the special fibre, but has a pole on $C$ and does not extend to $Z$.

For a minimal elliptic model, Corollary 3.27 gives $\omega_{X/S}=\pi^*M$ for the invertible $M=\pi_*\omega_{X/S}$. Trivialize $M$ locally on $S$. A section on $U$ is then a regular function times that frame, and (a), applied over that open of $S$, extends it uniquely. Gluing proves the intended equality under the corrected minimality hypothesis. \hfill$\square$

\exercise{3.10}{Move a divisor away from finitely many points and reducible fibres, split it into horizontal and vertical parts, and use torsion of $\operatorname{Pic}(S)$.}
\pf The set $S_0$ of reducible fibres is finite by generic geometric irreducibility and openness of geometric irreducibility in this flat family. Apply Proposition 1.11 to the $x_i$ and the finitely many generic points of components above $S_0$. This gives $D'\sim D$ avoiding them all. Decompose its Weil divisor into horizontal and vertical primes. Regularity makes both parts Cartier, giving $D'=H+V$. Each fibre contributing to $V$ is irreducible, say $X_s=d_s\Gamma_s$, so
\[ V=\sum_s n_s\Gamma_s=\sum_s\frac{n_s}{d_s}X_s. \]
Only finitely many terms occur. Clear their denominators to get $qV=\pi^*B$. If $\operatorname{Pic}(S)$ is torsion, a further positive multiple makes $B$ principal. After multiplying $q$ accordingly, $qD\sim qH$. The supports of $H$ and $V$ consist of disjoint sets of prime divisors, so their union is $\Supp D'$; in particular $H$ avoids all the prescribed points. \hfill$\square$

\section*{9.4 Contractions, canonical models and Weierstrass models}

\exercise{4.1}{Prove termination and correctness of the component-addition algorithm for a fundamental cycle.}
\pf Let $Z$ be the least positive antinef cycle supplied by Lemma 4.12, and let $\Delta\le A\le Z$ be a stage. If $A\cdot\Gamma>0$, the coefficient of $\Gamma$ in $A$ must be smaller than in $Z$. Otherwise $(Z-A)\cdot\Gamma\ge0$ because its $\Gamma$ coefficient is zero, contrary to $Z\cdot\Gamma\le0<A\cdot\Gamma$. Thus $A+\Gamma\le Z$. Every step increases the sum of coefficients, which is bounded by the corresponding sum for $Z$. At termination $A$ is antinef and contains $\Delta$, so minimality gives $Z\le A$ and hence $A=Z$. If $\Delta$ is already antinef, the algorithm has zero steps; the print's $n\ge1$ should allow this case. \hfill$\square$

\exercise{4.2}{Bound the fundamental cycle of a proper negative-definite fibre configuration by the full fibre.}
\pf Write $X_s=Z_1+Z_2$, where $Z_1$ has exactly the prime components in $\Delta$ and their original multiplicities. For a component $\Gamma$ of $\Delta$,
\[ Z_1\cdot\Gamma=-Z_2\cdot\Gamma\le0. \]
Thus $Z_1$ is an antinef cycle containing $\Delta$, so $Z\le Z_1$. Negative definiteness excludes an entire connected fibre component from $\Delta$, since its full fibre divisor would have square zero. Hence $Z_2>0$ and $Z\le Z_1<X_s$. The inequality means inequality of effective divisors, not strict inequality in every coefficient. \hfill$\square$

\exercise{4.3}{Over an algebraically closed residue field, compare $H^1(Z,\mathcal O_Z)=0$ and triviality of the multidegree-zero Picard group; analyze the reduced and nilpotent induction steps.}
\pf (a) For an integral proper curve, the connected degree-zero Picard group has dimension $h^1(\mathcal O_Z)$. Its abelian part is the Jacobian of the normalization, and its affine part comes from units at the singularities. Over an algebraically closed field a positive-dimensional connected algebraic group has nonidentity points. Thus it has trivial group of points exactly when its dimension is zero, equivalently $H^1=0$.

(b) Put $T=Z_1\cap Z_2$. The correct additive and multiplicative sequences are
\[ 0\to\mathcal O_Z\to\mathcal O_{Z_1}\oplus\mathcal O_{Z_2}
\to\mathcal O_T\to0, \]
\[ 1\to\mathcal O_Z^*\to\mathcal O_{Z_1}^*\times\mathcal O_{Z_2}^*
\to\mathcal O_T^*\to1. \]
The last maps are respectively difference and quotient. In particular
\[ 0\to\operatorname{coker}\bigl(H^0\mathcal O_{Z_1}\oplus H^0\mathcal O_{Z_2}
\to H^0\mathcal O_T\bigr)\to H^1\mathcal O_Z
\to H^1\mathcal O_{Z_1}\oplus H^1\mathcal O_{Z_2}\to0, \]
and
\[ \operatorname{coker}\bigl(H^0\mathcal O_{Z_1}^*\times H^0\mathcal O_{Z_2}^*
\to H^0\mathcal O_T^*\bigr)\hookrightarrow\operatorname{Pic}(Z)
\twoheadrightarrow\operatorname{Pic}(Z_1)\oplus\operatorname{Pic}(Z_2). \]
\textbf{The printed direct sums omit these cokernels.} For example take two smooth $(1,1)$ curves in $\mathbf P^1\times\mathbf P^1$ meeting transversally at two points. Each is $\mathbf P^1$, but their union has $H^1=k$ and $\operatorname{Pic}^0=k^*$, while both component groups vanish. The configuration is a valid vertical divisor for a fibred surface over $\Spec k$. For a general reduced curve, the additive gluing dimension and the dimension of the multiplicative gluing group are equal; both are the length of $T$ minus the rank of the map on constants. Together with (a), this proves the required equivalence for reduced $Z$, using the corrected exact sequences.

(c) Put $W=Z-\Gamma$ when $2\Gamma\le Z$. Its ideal in $Z$ is square-zero and equals $J=\mathcal O_\Gamma(-W)$. Hence
\[ 0\to J\to\mathcal O_Z\to\mathcal O_W\to0,\qquad
1\to1+J\to\mathcal O_Z^*\to\mathcal O_W^*\to1,\quad 1+J\simeq(J,+). \]
Both restriction maps on $H^1$ and on Picard groups are surjective, since $H^2(J)=0$. The additive kernel is $H^1(J)$ modulo the boundary of global functions; the multiplicative kernel is the same vector group modulo the boundary of global units. The following useful form of the units calculation explains why their vanishing is equivalent, even in mixed characteristic. Filter the finite rings of global functions by their nilpotent ideals. On each positive graded layer, $1+a\mapsto a$ identifies multiplication with addition; on the residue layer a nonzero scalar unit lifts locally and has zero additive boundary. Thus the dimension of the image of the units boundary equals the dimension of the image of the functions boundary. The multiplicative quotient is a connected unipotent group of dimension
\[ \dim_k\operatorname{coker}\bigl(H^0(W,\mathcal O_W)\to H^1(\Gamma,J)\bigr). \]
Here the boundary target is a $k$-vector space, and the dimension is the dimension of its additive quotient; nilpotent layers are taken successively if the source has no coefficient field. A connected unipotent quotient over algebraically closed $k$ is trivial exactly when this dimension is zero. This proves that $\operatorname{Pic}^0(Z)\to\operatorname{Pic}^0(W)$ is an isomorphism exactly when $H^1(Z,\mathcal O_Z)\to H^1(W,\mathcal O_W)$ is. The argument uses square-zero ideals killed by the residue characteristic, not the assumption that $Z$ itself is a curve over $k$.

(d) Remove repeated components one at a time. Each $H^1$ restriction and each $\operatorname{Pic}^0$ restriction is surjective. By (c) its kernel vanishes on one side exactly when it vanishes on the other. The reduced case established using the corrected (b) therefore gives $H^1(Z,\mathcal O_Z)=0$ if and only if $\operatorname{Pic}^0(Z)=1$. Disconnected supports can be treated componentwise. \hfill$\square$

\exercise{4.4}{From $H^1(Z,\mathcal O_Z)=0$ for a fundamental cycle, deduce the vanishing and Picard assertions for all cycles supported there.}
\pf For $0<Z'\le Z$, the restriction $H^1(Z,\mathcal O_Z)\to H^1(Z',\mathcal O_{Z'})$ is surjective, because its kernel sheaf has support of dimension one and no $H^2$. Thus $H^1(Z',\mathcal O_{Z'})=0$. By 4.2, $Z'\le X_s$ in the arithmetic case; in the geometric case it is already a curve over the base field. Its genus is consequently $1-h^0(\mathcal O_{Z'})\le0$. Lemma 4.5 gives $\operatorname{Pic}^0(Z)=1$, over arbitrary residue fields as well.

Since $Z\cdot\Gamma\le0$, the invertible sheaves $\mathcal O_Z(-nZ)$ have nonnegative component degrees. Lemma 4.11 gives both generation and $H^1(Z,\mathcal O_Z(-nZ))=0$. The sequences
\[ 0\to\mathcal O_Z(-nZ)\to\mathcal O_{(n+1)Z}\to\mathcal O_{nZ}\to0 \]
prove $H^1(nZ,\mathcal O_{nZ})=0$ inductively. Their units sequences have zero $H^1$ kernel, so $\operatorname{Pic}^0((n+1)Z)\hookrightarrow\operatorname{Pic}^0(nZ)$; starting at $Z$ gives triviality for every $n$. Finally any $Z'$ on the same support satisfies $Z'\le nZ$ for large $n$. Restriction gives $H^1(Z',\mathcal O_{Z'})=0$. To prove Picard triviality without assuming $Z'$ is itself a curve over the residue field, extend a multidegree-zero line bundle from $Z'$ to $nZ$. First increase the multiplicities on its existing support. These are nilpotent thickenings; their line-bundle lifting obstructions lie in $H^2$ of coherent sheaves and vanish, and component degrees are unchanged. Then add the remaining components with the multiplicities in $nZ$, giving them the trivial line bundle. Their intersection with the existing part is finite, and every invertible sheaf on that Artinian scheme is free. Choose such a trivialization and glue. This constructs a multidegree-zero line bundle on $nZ$ restricting to the given one. Since $\operatorname{Pic}^0(nZ)=1$, the original line bundle is trivial. These are precisely the vanishing and triviality conditions of Theorem 4.15. \hfill$\square$

\exercise{4.5}{Contract any vertical configuration containing no connected fibre component when $p_a(X_\eta)\le0$.}
\pf Work on a connected component of a closed fibre, replacing $S$ by the local Stein base when necessary. For clarity suppose first that $X_s$ is connected and put $F=X_s/d$, its full primitive fibre. The valid full-fibre argument in 1.14 injects $H^0(F,\mathcal O_F)$ into the field of constants of $F_{\mathrm{red}}$. It is therefore itself a field, finite over $k(s)$. Since $F^2=0$, adjunction gives
\[ \chi_k(\mathcal O_F)=\frac{1-p_a(X_\eta)}d>0. \]
Cohomology is a vector space over $l=H^0(F,\mathcal O_F)$ and $h_l^0=1$. Positive Euler characteristic then forces $h_l^1=0$. Thus $H^1(F,\mathcal O_F)=0$. Lemma 4.5 gives $\operatorname{Pic}^0(F)=1$, and $\mathcal O_F(-nF)$ has multidegree zero. The sequences for successive multiples of $F$ now give $H^1(X_s,\mathcal O_{X_s})=0$.

If the generic constants are larger than $K(S)$, use the finite normal Stein base $S_1$. The generic curve over $K(S_1)$ has genus zero because $\chi_K(\mathcal O_{X_\eta})>0$. Apply the preceding argument at each point of $S_1$; the same sequences for multiples handle the ramification thickness of the original fibre. Thus every subcurve of the original special fibre has $H^1=0$.

For the configuration $\Delta$, its matrix is negative definite by 1.6. Its fundamental cycle $Z$ satisfies $Z<X_s$ by 4.2 and hence has $H^1=0$. Exercise 4.4 supplies all conditions of Theorem 4.15 (in particular those of Theorem 4.2), so $\Delta$ contracts. Treat its finitely many connected components separately. \hfill$\square$

\exercise{4.6}{Prove that a canonical fibre has at most $2p_a(X_\eta)-2$ geometric components.}
\pf Base change the fibre to $\overline{k(s)}$. Relative ampleness of $\omega_{Y/S}$ persists, so its degree on every geometric prime component is a positive integer. If their multiplicities are $m_j\ge1$, invariance of degree under field extension gives
\[ 2p_a(X_\eta)-2=\deg\omega_{Y_s/k(s)}
=\sum_jm_j\deg\bigl(\omega_{Y_s/k(s)}|_{C_j}\bigr)\ge\sum_jm_j\ge\#\{C_j\}. \]
The degree formula includes possible nonreduced geometric fibres, and so needs no geometric reducedness assumption. \hfill$\square$

\exercise{4.7}{For the rational contraction of Theorem 4.15, identify the pullback of the maximal ideal, factor through the normalized blow-up, examine the canonical resolution sequence, and compute the fibre tangent-space criterion.}
\pf (a) In the notation of Lemma 4.14, $J=\mathcal O_X(-Z)$ is generated over an affine neighborhood of $y$ by $H^0(J)=\fm_y$. Hence $\fm_y\mathcal O_X=\mathcal O_X(-Z)$.

(b) This invertibility factors $f$ through $\operatorname{Bl}_yY$. Normality of $X$ factors it further through its normalization $Y_1$. The pullback of the relatively ample tautological sheaf is $\mathcal O_X(-Z)$; thus exactly the components $\Gamma$ with $Z\cdot\Gamma=0$ contract to $Y_1$. Their intersection matrix is a principal submatrix of a negative-definite matrix. Every cycle on their support already has the vanishing and Picard properties in Theorem 4.15, so the new contraction satisfies its hypotheses.

(c) \textbf{A prescribed $X$ need not be the endpoint of that sequence.} Take a regular $Y$, blow up one closed point, then a point of the exceptional curve, and let $f:X\to Y$ be the composite. Its exceptional configuration satisfies Theorem 4.15: the curves form a negative-definite rational tree, and the blow-up calculation or 4.4 gives the required vanishings for all cycles. Yet the canonical sequence of normalization and blowing up singular points of $Y$ has length zero, ending at the already regular $Y$, not at $X$. A version starting with singular $Y$ is obtained by taking a rational double point and then blowing up an extra smooth point of its minimal resolution. The valid conclusion from (b) is that the algorithm is dominated by $X$ and ends with a regular surface dominated by $X$; the remaining map is a sequence of point blow-ups. Equality with $X$ requires an additional minimality condition and, for a particular resolution algorithm, compatibility with its centers.

(d) Let $t$ generate the base maximal ideal. The closed-fibre tangent space is the annihilator of the class of $t$ in $\fm_y/\fm_y^2$. Therefore $T_{Y_s,y}=T_{Y,y}$ exactly when $t\in\fm_y^2$. Lemma 4.14 gives, for every $n$, $f_*\mathcal O_X(-nZ)=\fm_y^n$ near $y$. Thus $t\in\fm_y^2$ exactly when its divisor on $X$ contains $2Z$, that is $X_s\ge2Z$. Both directions use the equality of ideals, not just integral closure of an ideal. \hfill$\square$

\exercise{4.8}{For a fibre consisting of two elliptic curves meeting transversally at a rational point, test generation of powers of the relative canonical sheaf.}
\pf Localize at $s$, with residue field $k$, and call the components $E_1,E_2$ and their node $p$. The fibre is reduced with arithmetic genus two. Adjunction on this nodal curve gives
\[ \omega_{X/S}|_{E_i}=\omega_{E_i/k}(p)=\mathcal O_{E_i}(p). \]
Its one-dimensional space of sections vanishes at $p$. Thus every section of the fibre's dualizing sheaf, and a fortiori every section on $X$, vanishes in its fibre at $p$. The canonical sheaf is not generated there.

For its square, the restrictions are $\mathcal O_{E_i}(2p)$, base-point-free degree-two sheaves. Their evaluation maps at $p$ are surjective and their $H^1$ vanish. The gluing exact sequence at the node therefore shows that $\omega_{X_s}^{\otimes2}$ is globally generated and has $H^1=0$. Cohomology and base change lift these sections over the local base; Nakayama gives relative generation in a neighborhood of the special fibre. Properness over a local base makes this all of $X$. The generic fibre also has genus two and its bicanonical sheaf is generated (degree $4=2g$). Hence over $\Spec\mathcal O_{S,s}$ the least positive integer is $n=2$.

\textbf{No such global $n$ follows over an unrestricted $S$.} Blow up a smooth rational point of a different fibre of any family having the indicated fibre. On the new exceptional curve $E$, every $\omega_{X/S}^{\otimes n}$ has degree $-n$, so no positive power is generated, although the specified fibre has not changed. Even relative generation must be distinguished from generation by global sections over a nonaffine base. Thus two is the local answer; the global question needs additional hypotheses. \hfill$\square$

\exercise{4.9}{Characterize residue extensions at points of $\mathbf P^1_k$ and construct the integral genus-zero fibre with normalization $\mathbf P^1_{k'}$; show failure of l.c.i. for extension degree at least three.}
\pf (a) A point off infinity has residue field $k[T]/(f)$ for an irreducible polynomial, hence a simple extension. Infinity has residue field $k$, also simple. Conversely a primitive generator gives its minimal polynomial and the required point.

(b) Blow up that point in $\mathbf P^1_R$. The special fibre is $\Gamma+E$, with $E=\mathbf P^1_{k'}$, $\Gamma=\mathbf P^1_k$, and $\Gamma^2=-m$, $m=[k':k]$. The negative curve $\Gamma$ contracts by Theorem 4.15: its successive ideal quotients are $\mathcal O_{\mathbf P^1}(jm)$, all with $H^1=0$, and the degree-zero Picard groups of its thickenings vanish by the units sequence. Let $Y$ be the resulting normal surface. Its special fibre $C$ is integral, being the image of $E$ with multiplicity one. The finite birational map $E\to C$ is its normalization.

The rational contraction has $f_*\mathcal O=\mathcal O_Y$ and $R^1f_*\mathcal O=0$, by Lemma 4.14. Pushing the multiplication-by-$t$ sequence down shows the same statements on the special fibres. Since the blow-up of $\mathbf P^1_R$ has $H^1$ of its special fibre zero, so does $C$. Its constants are $k$: the contracted point has residue field $k$, which already forces the constant field of the integral proper curve $C$ to be $k$. Thus $p_a(C)=0$.

(c) If $C$ were l.c.i., its dualizing sheaf would be invertible of degree $-2$. Its pullback to the normalization has the same degree over $k$, since normalization is finite birational. Every invertible sheaf on $\mathbf P^1_{k'}$ has degree over $k$ divisible by $m$. Consequently $m\mid2$, impossible for $m\ge3$. \hfill$\square$

\exercise{4.10}{For an integral projective genus-one curve with a smooth rational point $p$, prove that $\mathcal O_C(3p)$ embeds it as a plane cubic, including the singular-case equations.}
\pf (a) The normalization $C'\to C$ has a rational point above the smooth $p$, so its constant field is $k$. Smoothness at $p$ also makes its function field separably generated over $k$. Thus that field is a regular extension of $k$, and $C$ is geometrically integral. One can check the last passage on affine charts: their coordinate domains embed into the function field, and remain domains after any field extension.

(b) Very ampleness and a closed immersion descend faithfully flatly, so extend $k$ to an algebraic closure. In the smooth case Riemann--Roch gives $h^0(3p)=3$ and separates length-two subschemes, since a degree-one line bundle has $H^1=0$ on a smooth genus-one curve. Hence it is very ample. It remains to handle the singular case.

(c) The normalization sequence gives $1=p_a(C)=g(C')+\delta$, with $\delta>0$. Hence $C'=\mathbf P^1_k$ and $\delta=1$. Take its coordinate $T$ with pole at the point above $p$. The affine ring $B=\mathcal O_C(C\setminus\{p\})$ has codimension one in $k[T]$. Since $1\in B$ and $T\notin B$ (otherwise $B=k[T]$), we have $k[T]=B\oplus kT$.

(d) Subtracting the $T$-components of $T^2,T^3$ gives $x=T^2+aT$ and $y=T^3+bT$ in $B$. Their leading degrees two and three generate every integer at least two. Successively cancelling leading terms of any polynomial in $B$ by monomials in $x,y$ leaves a constant plus a multiple of $T$; the latter coefficient must vanish. Therefore $B=k[x,y]$. Direct substitution yields
\[ y^2+(3ax+a^3+ab)y=x^3+2bx^2+(b^2+a^2b)x. \]
This identity is valid in every characteristic.

(e) The pole orders of $1,x,y$ at $p$ are zero, two and three. They form a basis of $H^0(C,\mathcal O(3p))$. On the affine complement they identify the image with $\Spec B$, the affine cubic just displayed. At infinity the projective cubic has the single smooth point $[0:1:0]$; $x/y$ is a uniformizer there and at $p$. The two affine identifications therefore glue to an isomorphism from $C$ onto that projective cubic. In particular the morphism is a closed immersion and $\mathcal O(3p)$ is very ample, including when $C$ is singular. \hfill$\square$

\exercise{4.11}{Extend the elliptic involution to an integral Weierstrass model and identify its quotient.}
\pf Completing the quadratic equation in the $y$-variable in its characteristic-free form shows that its two roots sum to $-a_1x-a_3$. Thus the involution is
\[ [u:v:z]\longmapsto[u:-v-a_1u-a_3z:z]. \]
It preserves the homogeneous equation over $A$ and fixes the origin. The rational coordinate $x=u/z$ has pole divisor $2o$ and defines a finite degree-two morphism $W\to\mathbf P^1_S$, including at the origin. Intrinsically it is given by the two generators $1,x$ of $\pi_*\mathcal O_W(2o)$; this description removes the apparent indeterminacy of the pair $[u:z]$ at $[0:1:0]$.

On $z\ne0$ the coordinate algebra is free with basis $1,y$ over $A[x]$. An invariant $b_0+b_1y$ must satisfy $2b_1=0$ and $b_1(a_1x+a_3)=0$. Since $A$ is a domain and the generic elliptic double cover is separable, these imply $b_1=0$, also in characteristic two. Hence the invariant algebra is $A[x]$. The same computation on the other base chart, or normality of $\mathbf P^1_S$ and finiteness of the quotient, gives the quotient everywhere. Thus $W/\langle\sigma\rangle=\mathbf P^1_S$. \hfill$\square$

\exercise{4.12}{Describe an isomorphism between integral Weierstrass equations of the pointed elliptic curve.}
\pf An isomorphism compatible with the pointed generic curve preserves the origin section and its pole-order filtration. The functions of pole order at most two have basis $1,x$ over $A$, and those of pole order at most three have basis $1,x,y$. The same holds for the primed model. Consequently
\[ x'=\alpha x+b,\qquad y'=\beta y+cx+d,
\quad\alpha,\beta\in A^*,\quad b,c,d\in A. \]
Comparing the leading pole-order-six terms in the equations gives $\beta^2=\alpha^3$. Set $a=\beta/\alpha\in A^*$. Then $a^2=\alpha$ and $a^3=\beta$, yielding $x'=a^2x+b$, $y'=a^3y+cx+d$. These are changes of equations of the \emph{pointed} curve, as usual for Weierstrass models. The repeated $a_i$ in the right side of the second printed equation are typographical; they should be the corresponding $b_i$. \hfill$\square$

\exercise{4.13}{Show that the local minimal discriminant does not depend on the chosen rational origin.}
\pf Translation by $o'-o$ identifies the two pointed generic elliptic curves. The uniqueness of the minimal regular model extends this generic isomorphism. It carries the closure of $o$ onto the closure of $o'$, hence carries the contraction in Theorem 4.35 to the corresponding contraction for the other origin. Over each local DVR it therefore identifies the two minimal Weierstrass models. Their equations differ by the integral change of variables of 4.12, whose scale is a unit. Discriminants change by its twelfth power, so their valuations are equal. \hfill$\square$

\exercise{4.14}{Identify the minimal-discriminant line bundle and deduce a global minimal equation under the stated torsion hypothesis.}
\pf Put $\lambda=\rho_*\omega_{X/S}$, an invertible sheaf. Choose local minimal equations with invariant differentials $\eta_i$ forming frames of $\lambda$. On overlaps a coordinate scale $a$ gives $\eta_j=a^{-1}\eta_i$ and $\Delta_j=a^{12}\Delta_i$. Thus $\Delta_i\eta_i^{12}$ is a global rational section of $\lambda^{12}$ with divisor $D$. Consequently $\mathcal O_S(D)\simeq\lambda^{12}$.

If $D$ is principal, $\lambda^{12}$ is trivial. The class of $\lambda^2$ has order dividing six, so the hypothesis makes it trivial. Then $\lambda$ has order dividing two and is trivial as well. Theorem 4.35 gives the criterion: over the affine Dedekind base a global minimal Weierstrass equation exists exactly when $\lambda$ is free. To see the coordinate content, the pole-order quotients for $1,x,y$ are the appropriate powers of $\lambda$; when it is free, their exact sequences of projective $A$-modules split. Global choices of $x,y$ give an integral equation, and the chosen invariant differential generates $\lambda$ at every prime, so the equation is minimal at every prime. This proves (b). \hfill$\square$

\exercise{4.15}{Show existence of a global minimal equation if every bad prime of a given integral equation is principal.}
\pf Let $\eta_0$ be the invariant differential of the given equation. At every good prime the equation is smooth and already minimal, so $\eta_0$ is a basis of the Hodge line $\lambda$ there. Thus as a fractional $A$-module inside $K\eta_0$,
\[ \lambda=\left(\prod_{\text{bad }\fp}\fp^{n_\fp}\right)\eta_0 \]
for integers $n_\fp$, with finitely many nonzero terms. Every ideal in this product is principal by hypothesis, including its inverse powers. Hence $\lambda$ is free. The global criterion in 4.14, equivalently Theorem 4.35, produces a global minimal Weierstrass equation. No principality assumption on the good primes is needed. \hfill$\square$

\exercise{4.16}{For a minimal genus-one arithmetic surface, analyze the effective vertical canonical divisor, its coefficient modulo a complete fibre, and triviality in the presence of a rational point.}
\pf The smooth geometrically connected generic genus-one curve has a one-dimensional space of differentials and canonical degree zero. Any nonzero differential has no zeros, so $\omega_{X_\eta/K}=\mathcal O_{X_\eta}$. Extending it rationally yields an integral vertical divisor $V$ with $\omega_{X/S}=\mathcal O(V)$. Nefness and $K\cdot X_s=0$ give $V\cdot\Gamma=0$ for every vertical prime. Hence on a connected fibre
\[ V_s=\frac{a_s}{d_s}X_s,\qquad a_s\in\Z,\quad
 d_s=\gcd\{\text{component multiplicities}\}. \]
Over a local DVR, write $a_s=q_sd_s+r_s$, $0\le r_s<d_s$, and subtract the principal divisor $q_sX_s=\operatorname{div}(t^{q_s})$. The resulting representative is the effective divisor $(r_s/d_s)X_s$. This proves (a),(b) \emph{locally}; the rational multiple is an integral divisor because every multiplicity is divisible by $d_s$.

If the generic fibre has a rational point, its closure is a section. Intersection with the fibre gives one, forcing $d_s=1$. Locally the preceding representative is zero, so $\omega_{X/S}$ is trivial over a local base. Over a general Dedekind base, the correct global conclusion is
\[ \omega_{X/S}\simeq\pi^*\lambda,\qquad
\lambda=\pi_*\omega_{X/S}\in\operatorname{Pic}(S), \]
as in Corollary 3.27; it is trivial if and only if $\lambda$ is trivial, since pulling back by the section recovers $\lambda$.

\textbf{Global effectivity and triviality in the printed assertions are false.} Let $S$ be a smooth projective curve over an algebraically closed field of characteristic zero with a nontrivial two-torsion line bundle $N$. Its associated connected etale double cover $T\to S$ has sign eigensheaf $N$. For a constant elliptic curve $E$, take the quotient of $E\times T$ by the diagonal action of inversion on $E$ and the deck involution on $T$. This is a smooth projective elliptic scheme $X\to S$, with its zero section; it is minimal because its geometric fibres are smooth genus-one curves. Its invariant differentials transform by the sign character, so $\lambda=N$ and $\omega_{X/S}=\pi^*N\ne\mathcal O_X$. Moreover $H^0(X,\omega)=H^0(S,N)=0$, since a nontrivial degree-zero line bundle on a proper curve has no section. An effective canonical representative would supply a nonzero section, so even global effectivity fails. The same triviality obstruction can occur over affine Dedekind bases. This distinguishes the valid local normalization from an unjustified global choice of uniformizers. \hfill$\square$

\end{document}
